\documentclass[11pt]{article}
\usepackage[margin=1in]{geometry}
\usepackage{amsmath,amssymb,mathtools}
\usepackage{enumitem}
\usepackage{hyperref}

\setlist[enumerate,1]{label=(\roman*), leftmargin=*}

\newcommand{\R}{\mathbb{R}}
\newcommand{\ii}{\iint}
\newcommand{\iii}{\iiint}

\title{Multivariable Calculus\\Problem Set 9: Solutions}
\author{Angad Singh Ahuja}
\date{\today}

\begin{document}
\maketitle

\section*{Problem 1}
Change of variable in multiple integrals.

\begin{enumerate}
\item Use $x=u^2-v^2$, $y=2uv$ to evaluate $\displaystyle \iint_R y\,dA$, where $R$ is bounded by the $x$-axis and the parabolas $y^2=4-4x$, $y^2=4+4x$, $y\ge 0$.

\textbf{Solution.}
Under the map
\[
x=u^2-v^2,\qquad y=2uv,
\]
the curves $u=1$ and $v=1$ become
\[
x=1-v^2,\ y=2v \quad\Longrightarrow\quad y^2=4-4x,
\]
\[
x=u^2-1,\ y=2u \quad\Longrightarrow\quad y^2=4+4x.
\]
The $x$-axis with $-1\le x\le 1$ corresponds to $u=0$ or $v=0$. Hence the given region is the image of the square
\[
0\le u\le 1,\qquad 0\le v\le 1.
\]
The Jacobian is
\[
J=\frac{\partial(x,y)}{\partial(u,v)}
=\begin{vmatrix}2u & -2v\\ 2v & 2u\end{vmatrix}
=4u^2+4v^2=4(u^2+v^2).
\]
Also $y=2uv$. Therefore
\[
\iint_R y\,dA
=\int_0^1\int_0^1 (2uv)\cdot 4(u^2+v^2)\,du\,dv.
\]
So
\[
=8\int_0^1\int_0^1 (u^3v+uv^3)\,du\,dv
=8\left(\frac14\cdot\frac12+\frac12\cdot\frac14\right)=8\cdot\frac14=2.
\]
Hence
\[
\boxed{ \iint_R y\,dA =2 }.
\]

\item Use $x=2u+v$, $y=u+2v$ to evaluate $\displaystyle \iint_R (x-3y)\,dA$, where $R$ has vertices $(0,0)$, $(2,1)$, $(1,2)$.

\textbf{Solution.}
The inverse linear system is not needed explicitly. We identify the corresponding $uv$-region from the vertices:
\[
(u,v)=(0,0)\mapsto (0,0),\qquad (1,0)\mapsto (2,1),\qquad (0,1)\mapsto (1,2).
\]
So $R$ is the image of the triangle
\[
T=\{(u,v):u\ge 0,\ v\ge 0,\ u+v\le 1\}.
\]
The Jacobian is
\[
J=\frac{\partial(x,y)}{\partial(u,v)}
=\begin{vmatrix}2&1\\1&2\end{vmatrix}=3.
\]
Also
\[
x-3y=(2u+v)-3(u+2v)=-u-5v.
\]
Hence
\[
\iint_R (x-3y)\,dA
=\int_0^1\int_0^{1-u} 3(-u-5v)\,dv\,du.
\]
Compute the inner integral:
\[
\int_0^{1-u}(-u-5v)\,dv=-u(1-u)-\frac52(1-u)^2.
\]
Therefore
\[
\iint_R (x-3y)\,dA
=3\int_0^1\left[-u(1-u)-\frac52(1-u)^2\right]du.
\]
Expanding,
\[
-u+u^2-\frac52+5u-\frac52u^2
=4u-\frac32u^2-\frac52.
\]
So
\[
3\int_0^1\left(4u-\frac32u^2-\frac52\right)du
=3\left[2u^2-\frac12u^3-\frac52u\right]_0^1
=3(2-\tfrac12-\tfrac52)=3(-1)=-3.
\]
Hence
\[
\boxed{\iint_R (x-3y)\,dA=-3 }.
\]

\item Use $x=2u$, $y=3v$ to evaluate $\displaystyle \iint_R x^2\,dA$, where $R$ is bounded by the ellipse $9x^2+4y^2=36$.

\textbf{Solution.}
Substituting $x=2u$, $y=3v$ into the ellipse,
\[
9(2u)^2+4(3v)^2=36
\quad\Longrightarrow\quad
36u^2+36v^2=36
\quad\Longrightarrow\quad
u^2+v^2\le 1.
\]
So the ellipse maps to the unit disk $u^2+v^2\le 1$.
The Jacobian is
\[
J=\frac{\partial(x,y)}{\partial(u,v)}
=\begin{vmatrix}2&0\\0&3\end{vmatrix}=6.
\]
Since $x^2=(2u)^2=4u^2$,
\[
\iint_R x^2\,dA=\iint_{u^2+v^2\le 1} 4u^2\cdot 6\,du\,dv
=24\iint_{u^2+v^2\le 1} u^2\,du\,dv.
\]
Using polar coordinates in the $uv$-plane,
\[
u=r\cos\theta,\quad v=r\sin\theta,\quad 0\le r\le 1,\ 0\le \theta\le 2\pi,
\]
we get
\[
24\int_0^{2\pi}\int_0^1 r^2\cos^2\theta\cdot r\,dr\,d\theta
=24\left(\int_0^1 r^3\,dr\right)\left(\int_0^{2\pi}\cos^2\theta\,d\theta\right).
\]
Thus
\[
=24\cdot \frac14\cdot \pi=6\pi.
\]
Hence
\[
\boxed{\iint_R x^2\,dA=6\pi }.
\]
\end{enumerate}

\section*{Problem 2}
Evaluating triple integrals with cylindrical coordinates.

\begin{enumerate}
\item A solid lies within the cylinder $x^2+y^2=1$, below the plane $z=4$, and above the paraboloid $z=1-x^2-y^2$. Find its volume.

\textbf{Solution.}
In cylindrical coordinates, $r\le 1$, $0\le \theta\le 2\pi$, and
\[
1-r^2\le z\le 4.
\]
Hence
\[
V=\int_0^{2\pi}\int_0^1\int_{1-r^2}^{4} r\,dz\,dr\,d\theta.
\]
Integrating in $z$,
\[
V=\int_0^{2\pi}\int_0^1 (4-(1-r^2))r\,dr\,d\theta
=\int_0^{2\pi}\int_0^1 (3+r^2)r\,dr\,d\theta.
\]
So
\[
V=2\pi\int_0^1(3r+r^3)\,dr
=2\pi\left[\frac{3r^2}{2}+\frac{r^4}{4}\right]_0^1
=2\pi\left(\frac32+\frac14\right)=\frac{7\pi}{2}.
\]
Thus
\[
\boxed{V=\frac{7\pi}{2}}.
\]

\item Evaluate
\[
\int_{-2}^{2}\int_{-\sqrt{4-x^2}}^{\sqrt{4-x^2}}\int_{\sqrt{x^2+y^2}}^{2}(x^2+y^2)\,dz\,dy\,dx.
\]

\textbf{Solution.}
The projection in the $xy$-plane is the disk $x^2+y^2\le 4$, i.e. $0\le r\le 2$, $0\le \theta\le 2\pi$. Also
\[
z \text{ runs from } r \text{ to } 2.
\]
Since $x^2+y^2=r^2$, the integral becomes
\[
\int_0^{2\pi}\int_0^2\int_r^2 r^2\cdot r\,dz\,dr\,d\theta
=\int_0^{2\pi}\int_0^2\int_r^2 r^3\,dz\,dr\,d\theta.
\]
Therefore
\[
=2\pi\int_0^2 r^3(2-r)\,dr
=2\pi\int_0^2 (2r^3-r^4)\,dr.
\]
Hence
\[
=2\pi\left[\frac{r^4}{2}-\frac{r^5}{5}\right]_0^2
=2\pi\left(8-\frac{32}{5}\right)=\frac{16\pi}{5}.
\]
So
\[
\boxed{\frac{16\pi}{5}}.
\]
\end{enumerate}

\section*{Problem 3}
Evaluate the integrals by changing to cylindrical coordinates.

\begin{enumerate}
\item
\[
\int_{-2}^{2}\int_{-\sqrt{4-y^2}}^{\sqrt{4-y^2}}\int_{\sqrt{x^2+y^2}}^{2} xz\,dz\,dx\,dy.
\]

\textbf{Solution.}
The region is the same as in Problem 2(ii): $0\le r\le 2$, $0\le \theta\le 2\pi$, $r\le z\le 2$.
Now
\[
x=r\cos\theta,
\]
so the integrand becomes
\[
xz\,dV=(r\cos\theta)z\cdot r\,dz\,dr\,d\theta=r^2 z\cos\theta\,dz\,dr\,d\theta.
\]
Thus
\[
\int_0^{2\pi}\int_0^2\int_r^2 r^2z\cos\theta\,dz\,dr\,d\theta.
\]
Since
\[
\int_0^{2\pi}\cos\theta\,d\theta=0,
\]
the whole integral is zero:
\[
\boxed{0}.
\]

\item
\[
\int_{-3}^{3}\int_{0}^{\sqrt{9-x^2}}\int_0^{\sqrt{9-x^2-y^2}}\sqrt{x^2+y^2}\,dz\,dy\,dx.
\]

\textbf{Solution.}
The $xy$-projection is the upper semicircle $x^2+y^2\le 9$, $y\ge 0$, so in cylindrical coordinates
\[
0\le r\le 3,\qquad 0\le \theta\le \pi,\qquad 0\le z\le \sqrt{9-r^2}.
\]
Also $\sqrt{x^2+y^2}=r$, so the integral becomes
\[
\int_0^\pi\int_0^3\int_0^{\sqrt{9-r^2}} r\cdot r\,dz\,dr\,d\theta
=\int_0^\pi\int_0^3 r^2\sqrt{9-r^2}\,dr\,d\theta.
\]
Hence
\[
=\pi\int_0^3 r^2\sqrt{9-r^2}\,dr.
\]
Set $r=3\sin\phi$, so $dr=3\cos\phi\,d\phi$, and when $r$ goes from $0$ to $3$, $\phi$ goes from $0$ to $\pi/2$. Then
\[
r^2=9\sin^2\phi,\qquad \sqrt{9-r^2}=3\cos\phi.
\]
So
\[
\pi\int_0^{\pi/2} 9\sin^2\phi\cdot 3\cos\phi\cdot 3\cos\phi\,d\phi
=81\pi\int_0^{\pi/2}\sin^2\phi\cos^2\phi\,d\phi.
\]
Using $\sin^2\phi\cos^2\phi=\frac14\sin^2(2\phi)$,
\[
\int_0^{\pi/2}\sin^2\phi\cos^2\phi\,d\phi=\frac{\pi}{16}.
\]
Therefore
\[
\boxed{\frac{81\pi^2}{16}}.
\]
\end{enumerate}

\section*{Problem 4}
Evaluate the integrals by changing to spherical coordinates.

\begin{enumerate}
\item
\[
\int_0^1\int_0^{\sqrt{1-x^2}}\int_{\sqrt{x^2+y^2}}^{\sqrt{2-x^2-y^2}} xy\,dz\,dy\,dx.
\]

\textbf{Solution.}
The lower surface is the cone $z=\sqrt{x^2+y^2}$, i.e. $z=r$, which in spherical coordinates means
\[
\rho\cos\phi=\rho\sin\phi \quad\Longrightarrow\quad \phi=\frac{\pi}{4}.
\]
The upper surface is the sphere
\[
x^2+y^2+z^2=2 \quad\Longrightarrow\quad \rho=\sqrt2.
\]
Also $x,y\ge 0$, so
\[
0\le \theta\le \frac{\pi}{2},\qquad 0\le \phi\le \frac{\pi}{4},\qquad 0\le \rho\le \sqrt2.
\]
Now
\[
x=\rho\sin\phi\cos\theta,\qquad y=\rho\sin\phi\sin\theta,
\]
so
\[
xy=\rho^2\sin^2\phi\cos\theta\sin\theta.
\]
Multiplying by the Jacobian $\rho^2\sin\phi$ gives
\[
\rho^4\sin^3\phi\cos\theta\sin\theta.
\]
Therefore
\[
\int_0^{\pi/2}\int_0^{\pi/4}\int_0^{\sqrt2}\rho^4\sin^3\phi\cos\theta\sin\theta\,d\rho\,d\phi\,d\theta.
\]
Separate the integrals:
\[
\left(\int_0^{\sqrt2}\rho^4\,d\rho\right)
\left(\int_0^{\pi/4}\sin^3\phi\,d\phi\right)
\left(\int_0^{\pi/2}\cos\theta\sin\theta\,d\theta\right).
\]
These are
\[
\int_0^{\sqrt2}\rho^4\,d\rho=\frac{(\sqrt2)^5}{5}=\frac{4\sqrt2}{5},
\]
\[
\int_0^{\pi/2}\cos\theta\sin\theta\,d\theta=\frac12,
\]
\[
\int_0^{\pi/4}\sin^3\phi\,d\phi
=\left[-\cos\phi+\frac{\cos^3\phi}{3}\right]_0^{\pi/4}
=\frac23-\frac{5\sqrt2}{12}.
\]
Hence
\[
\boxed{\frac{4\sqrt2-5}{15}}.
\]

\item
\[
\int_{-a}^{a}\int_{-\sqrt{a^2-y^2}}^{\sqrt{a^2-y^2}}
\int_{-\sqrt{a^2-x^2-y^2}}^{\sqrt{a^2-x^2-y^2}}
(x^2z+y^2z+z^3)\,dz\,dx\,dy.
\]

\textbf{Solution.}
The region is the full ball $x^2+y^2+z^2\le a^2$, symmetric with respect to the plane $z=0$.
The integrand can be factored as
\[
x^2z+y^2z+z^3=z(x^2+y^2+z^2).
\]
This is an odd function of $z$. Since the region is symmetric in $z$, the integral is zero:
\[
\boxed{0}.
\]

\item
\[
\int_{-2}^{2}\int_{-\sqrt{4-x^2}}^{\sqrt{4-x^2}}
\int_{-\sqrt{4-x^2-y^2}}^{\sqrt{4-x^2-y^2}}
(x^2+y^2+z^2)^{3/2}\,dz\,dy\,dx.
\]

\textbf{Solution.}
The region is the full ball $\rho\le 2$.
Since $x^2+y^2+z^2=\rho^2$, the integrand becomes $\rho^3$.
Thus
\[
\int_0^{2\pi}\int_0^\pi\int_0^2 \rho^3\cdot \rho^2\sin\phi\,d\rho\,d\phi\,d\theta
=\int_0^{2\pi}d\theta\int_0^\pi \sin\phi\,d\phi\int_0^2 \rho^5\,d\rho.
\]
Hence
\[
=2\pi\cdot 2\cdot \frac{2^6}{6}
=4\pi\cdot \frac{64}{6}
=\boxed{\frac{128\pi}{3}}.
\]
\end{enumerate}

\section*{Problem 5}
Sketch the solid and evaluate the integral.

\begin{enumerate}
\item $\displaystyle \int_0^{\pi/6}\int_0^{\pi/2}\int_0^3 \rho^2\sin\phi\,d\rho\,d\theta\,d\phi$.

\textbf{Solution.}
This is the volume of the sector
\[
0\le \rho\le 3,\quad 0\le \theta\le \frac{\pi}{2},\quad 0\le \phi\le \frac{\pi}{6}.
\]
So
\[
V=\left(\int_0^3 \rho^2\,d\rho\right)\left(\int_0^{\pi/2}d\theta\right)\left(\int_0^{\pi/6}\sin\phi\,d\phi\right).
\]
Thus
\[
V=9\cdot \frac{\pi}{2}\cdot \left(1-\frac{\sqrt3}{2}\right)
=\boxed{\frac{9\pi}{4}(2-\sqrt3)}.
\]

\item $\displaystyle \int_0^{2\pi}\int_{\pi/2}^{\pi}\int_1^2 \rho^2\sin\phi\,d\rho\,d\phi\,d\theta$.

\textbf{Solution.}
This is the lower half of the spherical shell $1\le \rho\le 2$.
Hence
\[
V=\left(\int_1^2 \rho^2\,d\rho\right)\left(\int_{\pi/2}^{\pi}\sin\phi\,d\phi\right)\left(\int_0^{2\pi}d\theta\right).
\]
Compute:
\[
\int_1^2 \rho^2\,d\rho=\frac{8-1}{3}=\frac73,\qquad \int_{\pi/2}^{\pi}\sin\phi\,d\phi=1.
\]
Therefore
\[
V=\frac73\cdot 1\cdot 2\pi=\boxed{\frac{14\pi}{3}}.
\]

\item $\displaystyle \int_0^{\pi/2}\int_0^{\pi/2}\int_1^2 \rho^2\sin\phi\,d\rho\,d\phi\,d\theta$.

\textbf{Solution.}
This is the first-octant part of the shell $1\le \rho\le 2$.
So
\[
V=\left(\int_1^2 \rho^2\,d\rho\right)\left(\int_0^{\pi/2}\sin\phi\,d\phi\right)\left(\int_0^{\pi/2}d\theta\right)
=\frac73\cdot 1\cdot \frac{\pi}{2}
=\boxed{\frac{7\pi}{6}}.
\]

\item
\[
\int_{-2}^{2}\int_0^{\sqrt{4-y^2}}\int_{-\sqrt{4-x^2-y^2}}^{\sqrt{4-x^2-y^2}}
y^2\sqrt{x^2+y^2+z^2}\,dz\,dx\,dy.
\]

\textbf{Solution.}
The region is the half-ball $\rho\le 2$ with $x\ge 0$.
Thus
\[
0\le \rho\le 2,\quad 0\le \phi\le \pi,\quad -\frac{\pi}{2}\le \theta\le \frac{\pi}{2}.
\]
Now
\[
y=\rho\sin\phi\sin\theta,\qquad \sqrt{x^2+y^2+z^2}=\rho.
\]
So the integrand becomes
\[
y^2\sqrt{x^2+y^2+z^2}=\rho^3\sin^2\phi\sin^2\theta.
\]
Multiplying by the Jacobian gives
\[
\rho^5\sin^3\phi\sin^2\theta.
\]
Hence
\[
\int_{-\pi/2}^{\pi/2}\int_0^\pi\int_0^2 \rho^5\sin^3\phi\sin^2\theta\,d\rho\,d\phi\,d\theta.
\]
Separate:
\[
\left(\int_0^2 \rho^5\,d\rho\right)
\left(\int_0^\pi \sin^3\phi\,d\phi\right)
\left(\int_{-\pi/2}^{\pi/2}\sin^2\theta\,d\theta\right).
\]
These are
\[
\frac{64}{6}=\frac{32}{3},\qquad \frac{4}{3},\qquad \frac{\pi}{2}.
\]
Therefore
\[
\boxed{\frac{64\pi}{9}}.
\]
\end{enumerate}

\section*{Problem 6}
Solve the following.

\begin{enumerate}
\item Use spherical coordinates to find the volume of the solid above the cone $z=\sqrt{x^2+y^2}$ and below the sphere $x^2+y^2+z^2=z$.

\textbf{Solution.}
The sphere is
\[
\rho^2=\rho\cos\phi \quad\Longrightarrow\quad \rho=\cos\phi.
\]
The cone is $\phi=\pi/4$, and ``above the cone'' means $0\le \phi\le \pi/4$.
Therefore
\[
0\le \theta\le 2\pi,\qquad 0\le \phi\le \frac{\pi}{4},\qquad 0\le \rho\le \cos\phi.
\]
Thus
\[
V=\int_0^{2\pi}\int_0^{\pi/4}\int_0^{\cos\phi}\rho^2\sin\phi\,d\rho\,d\phi\,d\theta.
\]
Integrating in $\rho$,
\[
V=\frac{2\pi}{3}\int_0^{\pi/4}\cos^3\phi\sin\phi\,d\phi.
\]
Let $u=\cos\phi$, $du=-\sin\phi\,d\phi$:
\[
\int_0^{\pi/4}\cos^3\phi\sin\phi\,d\phi
=-\int_1^{\sqrt2/2}u^3\,du
=\frac{1}{4}\left(1-\frac14\right)=\frac{3}{16}.
\]
Hence
\[
\boxed{V=\frac{\pi}{8}}.
\]

\item Evaluate $\displaystyle \iiint_B e^{(x^2+y^2+z^2)^{3/2}}\,dV$, where $B$ is the unit ball.

\textbf{Solution.}
In spherical coordinates, $x^2+y^2+z^2=\rho^2$, so
\[
(x^2+y^2+z^2)^{3/2}=\rho^3.
\]
Since $B$ is the unit ball,
\[
0\le \rho\le 1,\qquad 0\le \phi\le \pi,\qquad 0\le \theta\le 2\pi.
\]
Therefore
\[
\iiint_B e^{(x^2+y^2+z^2)^{3/2}}\,dV
=\int_0^{2\pi}\int_0^\pi\int_0^1 e^{\rho^3}\rho^2\sin\phi\,d\rho\,d\phi\,d\theta.
\]
So
\[
=4\pi\int_0^1 e^{\rho^3}\rho^2\,d\rho.
\]
Let $u=\rho^3$, $du=3\rho^2\,d\rho$:
\[
4\pi\cdot \frac13\int_0^1 e^u\,du
=\frac{4\pi}{3}(e-1).
\]
Hence
\[
\boxed{\frac{4\pi}{3}(e-1)}.
\]

\item The plane $\frac{x}{a}+\frac{y}{b}+\frac{z}{c}=1$ cuts the ellipsoid $\frac{x^2}{a^2}+\frac{y^2}{b^2}+\frac{z^2}{c^2}\le 1$ into two pieces. Find the volume of the smaller piece.

\textbf{Solution.}
Use the linear change of variables
\[
x=au,\qquad y=bv,\qquad z=cw.
\]
Then the ellipsoid becomes the unit ball
\[
u^2+v^2+w^2\le 1,
\]
and the plane becomes
\[
u+v+w=1.
\]
The Jacobian is $abc$, so volumes in the ellipsoid are $abc$ times the corresponding volumes in the unit ball.

The plane $u+v+w=1$ has distance
\[
d=\frac{1}{\sqrt{1^2+1^2+1^2}}=\frac{1}{\sqrt3}
\]
from the origin. Hence the smaller piece is a spherical cap of the unit sphere of height
\[
h=1-\frac{1}{\sqrt3}.
\]
The volume of a cap of height $h$ in a sphere of radius $1$ is
\[
V_{\text{cap}}=\frac{\pi h^2(3-h)}{3}.
\]
Therefore the required volume is
\[
\boxed{abc\cdot \frac{\pi}{3}\left(1-\frac{1}{\sqrt3}\right)^2\left(2+\frac{1}{\sqrt3}\right)}.
\]

\item Find the volume of the solid within both the cylinder $x^2+y^2=1$ and the sphere $x^2+y^2+z^2=4$.

\textbf{Solution.}
In cylindrical coordinates, the common region is
\[
0\le r\le 1,\qquad 0\le \theta\le 2\pi,\qquad -\sqrt{4-r^2}\le z\le \sqrt{4-r^2}.
\]
Thus
\[
V=\int_0^{2\pi}\int_0^1\int_{-\sqrt{4-r^2}}^{\sqrt{4-r^2}} r\,dz\,dr\,d\theta
=2\pi\int_0^1 2r\sqrt{4-r^2}\,dr.
\]
So
\[
V=4\pi\int_0^1 r\sqrt{4-r^2}\,dr.
\]
Let $u=4-r^2$, $du=-2r\,dr$:
\[
V=4\pi\cdot \frac12\int_3^4 u^{1/2}\,du
=2\pi\cdot \frac23\left(4^{3/2}-3^{3/2}\right).
\]
Hence
\[
\boxed{\frac{4\pi}{3}(8-3\sqrt3)}.
\]
\end{enumerate}

\section*{Problem 7}
Determine whether or not $F$ is conservative in $\R^2$. If yes, find $\phi$ such that $F=\nabla\phi$.

\begin{enumerate}
\item $F=(2x-3y,\,-3x+4y-8)$.

\textbf{Solution.}
Let $P=2x-3y$, $Q=-3x+4y-8$. Then
\[
P_y=-3,\qquad Q_x=-3.
\]
So $F$ is conservative.
Find a potential by integrating $P$ with respect to $x$:
\[
\phi(x,y)=\int (2x-3y)\,dx=x^2-3xy+g(y).
\]
Now
\[
\phi_y=-3x+g'(y)=Q=-3x+4y-8,
\]
so
\[
g'(y)=4y-8 \quad\Longrightarrow\quad g(y)=2y^2-8y.
\]
Hence
\[
\boxed{\phi(x,y)=x^2-3xy+2y^2-8y +C }.
\]

\item $F=(e^x\sin y,\ e^x\cos y)$.

\textbf{Solution.}
\[
P_y=e^x\cos y,\qquad Q_x=e^x\cos y.
\]
So $F$ is conservative. Integrate $P$ in $x$:
\[
\phi(x,y)=\int e^x\sin y\,dx=e^x\sin y + g(y).
\]
Then
\[
\phi_y=e^x\cos y+g'(y)=Q=e^x\cos y,
\]
so $g'(y)=0$. Therefore
\[
\boxed{\phi(x,y)=e^x\sin y + C }.
\]

\item $F=(3x^2-2y^2,\ 4xy+3)$.

\textbf{Solution.}
\[
P_y=-4y,\qquad Q_x=4y.
\]
Since $P_y\ne Q_x$, $F$ is \emph{not} conservative.

\item $F=(ye^x+\sin y,\ e^x+x\cos y)$.

\textbf{Solution.}
\[
P_y=e^x+\cos y,\qquad Q_x=e^x+\cos y.
\]
So $F$ is conservative. Integrate $P$ in $x$:
\[
\phi(x,y)=\int (ye^x+\sin y)\,dx=ye^x+x\sin y+g(y).
\]
Then
\[
\phi_y=e^x+x\cos y+g'(y)=Q=e^x+x\cos y,
\]
so $g'(y)=0$. Hence
\[
\boxed{\phi(x,y)=ye^x+x\sin y + C }.
\]

\item $F=(2xy+y^{-2},\ x^2-2xy^{-3})$.

\textbf{Solution.}
Here the field is defined on $y\ne 0$. Compute
\[
P_y=2x-2y^{-3},\qquad Q_x=2x-2y^{-3}.
\]
So the field is conservative on each connected component of its domain.
Integrate $P$ in $x$:
\[
\phi(x,y)=\int (2xy+y^{-2})\,dx=x^2y+xy^{-2}+g(y).
\]
Then
\[
\phi_y=x^2-2xy^{-3}+g'(y)=Q=x^2-2xy^{-3},
\]
so $g'(y)=0$. Therefore
\[
\boxed{\phi(x,y)=x^2y+xy^{-2}+C }.
\]
\end{enumerate}

\section*{Problem 8}
Determine whether the vector field is conservative in $\R^3$. If yes, find $\phi$ such that $F=\nabla\phi$.

\begin{enumerate}
\item $F=(y^2z^3,\ 2xyz^3,\ 3xy^2z^2)$.

\textbf{Solution.}
Observe that
\[
\phi(x,y,z)=xy^2z^3
\]
has gradient
\[
\nabla\phi=(y^2z^3,\ 2xyz^3,\ 3xy^2z^2).
\]
Hence $F$ is conservative, with
\[
\boxed{\phi(x,y,z)=xy^2z^3 + C }.
\]

\item $F=(xyz^2,\ x^2yz^2,\ x^2y^2z)$.

\textbf{Solution.}
Take
\[
P=xyz^2,\qquad Q=x^2yz^2.
\]
Then
\[
P_y=xz^2,\qquad Q_x=2xyz^2.
\]
These are not equal in general, so $F$ is \emph{not} conservative.

\item $F=(3xy^2z^2,\ 2x^2yz^3,\ 3x^2y^2z^2)$.

\textbf{Solution.}
Again test cross-partials:
\[
P_y=6xyz^2,\qquad Q_x=4xyz^3.
\]
These are not equal in general, so $F$ is \emph{not} conservative.

\item $F=(e^{yz},\ xze^{yz},\ xye^{yz})$.

\textbf{Solution.}
Take
\[
\phi(x,y,z)=xe^{yz}.
\]
Then
\[
\phi_x=e^{yz},\qquad \phi_y=xze^{yz},\qquad \phi_z=xye^{yz}.
\]
Thus $F=\nabla\phi$, so
\[
\boxed{\phi(x,y,z)=xe^{yz}+C }.
\]

\item $F=(e^x\sin(yz),\ ze^x\cos(yz),\ ye^x\cos(yz))$.

\textbf{Solution.}
Take
\[
\phi(x,y,z)=e^x\sin(yz).
\]
Then
\[
\phi_x=e^x\sin(yz),\qquad \phi_y=ze^x\cos(yz),\qquad \phi_z=ye^x\cos(yz).
\]
Hence $F$ is conservative, with
\[
\boxed{\phi(x,y,z)=e^x\sin(yz)+C }.
\]
\end{enumerate}

\end{document}
